\documentclass[hyphens,aspectratio=169,dvipsnames]{beamer}
\usepackage{graphicx}
\usepackage{xcolor}
\usepackage{amssymb}
\usepackage{amsmath}
\usepackage{mathtools}
\usepackage{textcomp}
\usepackage{moresize}
\usepackage{framed}
\usepackage{minted}
\usepackage{relsize}
\usepackage{tikz}
\usepackage{comment}
\usetikzlibrary{shapes.geometric, arrows, positioning}
\usetheme{Berlin}
\usecolortheme[RGB={0,95,47}]{structure}
\beamertemplatenavigationsymbolsempty

\makeatletter
\AtBeginEnvironment{minted}{\dontdofcolorbox}
\def\dontdofcolorbox{\renewcommand\fcolorbox[4][]{##4}}
\makeatother

\newcommand{\textpf}[1]{\texttt{\color{black}\fcolorbox{lightgray}{lightgray}{#1}}}

\begin{document}

\begin{frame}{Learning objectives}

Today's learning objectives:
    \begin{itemize}
        \item Debugging with GDB
        \item Endianness
        \item How to make system calls in assembly
    \end{itemize}
    
\end{frame}

\begin{frame}
    \vspace{1cm}
    \hspace{5cm}
    \textbf{GDB}
\end{frame}


\begin{frame}{Debugging with GDB}
    \textit{``Debugging is twice as hard as writing a program in the first place. So if you're as clever as you can be when you write it, you are by definition not smart enough to debug it'' -- Brian Kernighan (Bell Labs, involved in the development of UNIX).}

\end{frame}

\begin{frame}{GDB}
    \textbf{Debugging capabilities}
    \begin{itemize}
        \item Pause program execution
        \item Inspect register/memory values
        \item Step through code instruction by instruction
        \item Edit register/memory values
    \end{itemize}

    \pause Here's a handout.
\end{frame}

\begin{frame}{Some useful commands}
    (More in CS 15 handout) \\ 
r: run \\ 
c: continue \\ 
b: breakpoint \\ 
s: stepi (step to next instruction) \\ 
\end{frame}

\begin{frame}
    \vspace{1cm}
    \hspace{5cm}
    \textbf{Endianness}
\end{frame}

\begin{frame}[fragile]{\textpf{mov}ing to and from Memory}
    \begin{itemize}
        \pause \item The \textpf{mov} instruction can read data from memory to registers, and from registers to memory.
        \pause \item To read 1 byte from memory at the address contained in \textpf{rsp} into the \textpf{al} register, you would use \textpf{mov al, byte ptr [rsp]}.
        \pause \item To read 2 bytes from memory at the address contained in \textpf{rsp} into the \textpf{ax} register, you would use \textpf{mov ax, word ptr [rsp]}.
        \pause \item To read 4 bytes from memory at the address contained in \textpf{rsp} into the \textpf{eax} register, you would use \textpf{mov eax, dword ptr [rsp]}.
        \pause \item To read 8 bytes from memory at the address contained in \textpf{rsp} into the \textpf{rax} register, you would use \textpf{mov rax, qword ptr [rsp]}.
    \end{itemize}
\end{frame}

\begin{frame}[fragile]{\textpf{.byte}}
    \begin{itemize}
        \pause \item \textpf{ret} is a one byte instruction. Its machine code is \textpf{0xc3}.
        \pause \item When you write \textpf{ret} in your program, the assembler writes a \textpf{0xc3} byte into its output machine code file.
        \pause \item You can also tell the assembler to emit this byte directly using the \textpf{.byte} directive:
        \pause \item \textpf{.byte 0xc3}.
        \pause \item Demo
    \end{itemize}
\end{frame}

\begin{frame}[fragile]{Data in \textpf{asm}}
    \begin{itemize}
        \pause \item Sometimes, you need to store more data than you can fit in the registers.
        \pause \item You can tell the assembler to use a portion of your program for data storage with the \textpf{.data} directive.
        \pause \item After that, use the \textpf{.byte} directive to reserve bytes provide an initial value for your variable.
        \pause \item When you're done, tell the assembler to switch back to ``code mode" using the \textpf{.text} directive.
        \pause \item Demo
    \end{itemize}
\end{frame}

\begin{frame}[fragile]{Motivating Example}
    What does this return?
    \begin{minted}{asm}
.data
some_cool_bytes:
    .byte 0xff, 0x00

.text
.global main
main:
    mov rax, 0
    mov ax, word ptr [some_cool_bytes]
    ret
    \end{minted}

    \pause A: It depends on the endianness of your CPU.
\end{frame}

\begin{frame}{Endianness}
    CPUs can be \textbf{little-endian} or \textbf{big-endian} (aka LSB/le/el, MSB/be/eb). \\

    When a little-endian CPU executes a \textbf{multi-byte} read, it reverses the order of the bytes before processing them. \\

    When a little-endian CPU executes a \textbf{multi-byte} write, it writes the bytes from least-significant to most-significant.
\end{frame}

\begin{frame}{Examples of Little-Endian Reads}
    Suppose that a computer's memory had the following contents at address \texttt{p}:
    \begin{center}
        \begin{tabular}{c|c|c|c|c|c}
            \hline
            ... & \texttt{0xDE} & \texttt{0xAD} & 0\texttt{xBE} & \texttt{0xEF} & ... \\
            \hline
        \end{tabular}

            $\overset{\uparrow}{\texttt{p}}~~~~~~~~~~~~~~~~~~~~~~~~~~~~~~~~~~~~~~$ \\ % This is a hack :)

    \end{center}

    When I read 2 bytes from \texttt{p} on a big-endian machine, I get... \pause \texttt{0xDEAD}.

    When I read 2 bytes from \texttt{p} on a little-endian machine, I get... \pause \texttt{0xADDE}.

    When I read 4 bytes from \texttt{p} on a big-endian machine, I get... \pause \texttt{0xDEADBEEF}.

    When I read 4 bytes from \texttt{p} on a little-endian machine, I get... \pause \texttt{0xEFBEADDE}.
\end{frame}

\begin{frame}{Endianness Misconception 1}
    Endianness is about bytes, \textbf{not bits}.

    A little-endian system reverses the order of bytes during memory accesses.

    \begin{center}\textbf{It does not change the order of the bits within each byte.}\end{center}

    As a consequence, if you always read and write things 1 byte at a time, you will not observe any differences between big- and little-endian systems.
\end{frame}

\begin{frame}[fragile]{Endianness Misconception 2}
    Endianness does not affect the direction of memory accesses.

    That is, if we execute a 4 byte read at the address \texttt{0x4000}, then we're reading from addresses \texttt{0x4000}, \texttt{0x4001}, \texttt{0x4002}, and \texttt{0x4003}, regardless of the endianness of the system.
\end{frame}

\begin{frame}[fragile]
    \vspace{1cm}
    \hspace{5cm}
    \textbf{ASM system calls}
\end{frame}

\begin{frame}[fragile]{Making System Calls in Assembly}
    On x86\_64, a system call is made like this:
    \begin{enumerate}
        \pause \item Place the system call number in \textpf{rax}.
        \pause \item Place the arguments in \textpf{rdi}, \textpf{rsi}, \textpf{rdx}, \textpf{r10}, \textpf{r8}, \textpf{r9}, in order.
        \pause \item Execute a \textpf{syscall} instruction.
        \pause \item (The OS executes the syscall)
        \pause \item The syscall returns, with the result left in \textpf{rax}. Additionally, \textpf{rcx} and \textpf{r11} are overwritten with junk values.
    \end{enumerate}
\end{frame}

\begin{frame}[fragile]{Activity: Your First Assembly System Call}
    Write an assembly program to exit with status 25, without returning from \textpf{main}.
\end{frame}

\begin{frame}[fragile]{Bypassing \textpf{main}: \textpf{\_start}}
    \begin{itemize}
        \pause \item We have previously mentioned that initialization code calls \textpf{main} for you, then makes an \textpf{exit} syscall with \textpf{main}'s returned value.
        \pause \item If you want to do that part yourself, define \textpf{\_start} instead of \textpf{main}.
        \pause \item \textpf{\_start} is the true entry point of the program.
        \pause \item When you do this, you need to tell the compiler driver not to provide its own \textpf{\_start}, so that yours will be used instead:
    \end{itemize}
    \pause \begin{center} \textpf{gcc -static -nostdlib your\_file.s} \end{center}
\end{frame}

\begin{frame}[fragile]{\textpf{argc} and \textpf{argv} in \textpf{\_start}}
    At the beginning of \textpf{\_start},
    \begin{itemize}
        \pause \item \textpf{argc} is written in memory as an 8-byte value at address \textpf{rsp}.
        \pause \item \textpf{argv} begins at address \textpf{rsp + 8}, and ends at \textpf{rsp + 8 * argc}
    \end{itemize}
\end{frame}

\begin{frame}[fragile]{Activity: Implement \textpf{echo} in Assembly, using \textpf{\_start} instead of \textpf{main}}
    \begin{itemize}
        \item \textpf{echo} is a Unix utility that writes each of its command-line arguments to \textpf{stdout}, separated by spaces, then exits 0.
        \item Implement \textpf{echo} in assembly.
    \end{itemize}
\end{frame}

\begin{frame}[fragile]{Activity: Implement \textpf{cat} in Assembly.}
    Much of this code will be reusable in the homework assignment that I will release today :)
\end{frame}

\end{document}
